\section{Introduction}
\label{sec:intro}
\IEEEPARstart{C}{hest} radiography is the most frequent imaging examination of the chest, but it projects the whole thorax onto one
plane: ribs, clavicles, heart and mediastinum overlap the lung fields and can hide small abnormalities. Chest CT resolves the same
anatomy in three dimensions, and it can therefore supervise radiograph models in ways that radiographs alone cannot. The bone and
soft-tissue components of a projected CT are known exactly, and findings can be learned from CT without the overlap of ribs, heart and
mediastinum. Both uses are now practical with public data, through projections rendered from chest CT \cite{angaitkar2026anatomy} and
through CT volumes paired with report-derived findings \cite{hamamci2026ctrate}.

CT-supervised models of radiographs are still judged mostly by image fidelity or classification accuracy. These aggregate measures hide
three clinically relevant failures. A model may \emph{erase} a lesion, so that a nodule present in the radiograph is missing from its
output; it may \emph{invent} a finding, reporting a disease that the image does not contain; or it may justify a correct output with a
heat map that points elsewhere. A suppressor with a high peak signal-to-noise ratio (PSNR) can still erase lesions, and a classifier with
a high area under the receiver operating characteristic curve (AUROC) can still respond to noise. This paper builds CT-supervised models
of chest radiographs and tests them for all three failures.

Dual-energy subtraction (DES) separates bone from soft tissue, but it needs dedicated hardware and a second exposure
\cite{kuhlman2006des,sun2026bsldm}. Software bone suppression avoids both. Radiologists find more nodules with suppressed images
\cite{freedman2011lung}, and suppressed images have been compared with computer-aided detection \cite{schalekamp2014cad} and used as
classifier inputs \cite{rajaraman2021bone}. Current learning-based suppressors either render paired supervision from chest CT, projecting
bone, non-lung soft tissue and lung separately \cite{angaitkar2026anatomy}, or train diffusion models on real DES pairs
\cite{chen2024bsdiff,sun2026bsldm,sun2025gllcm,sun2026debonedit}. In the CT-supervised line, lesion preservation has been assessed only
indirectly, through detection gains and the contrast of boxed lesions \cite{angaitkar2026anatomy}.

Suppressors are judged by how closely their output matches a bone or soft-tissue reference. A network rewarded for predicting bone can,
however, move lesion signal into the bone image, so that the lesion disappears from the suppressed radiograph. Because a model that does
nothing never absorbs a lesion, absorption has to be read next to suppression strength. Simulated nodules have been used to measure
contrast after hand-built clavicle suppression and in phantoms \cite{hogeweg2013suppression,takagi2020dose}, and learned image-to-image
models are known to remove or invent content in other modalities
\cite{cohen2018distribution,antun2020instabilities,bhadra2021hallucinations,tivnan2024hallucination}. In our search (PubMed, arXiv and
Crossref) we found no audit of learned bone suppression against an exactly known, additive lesion signal, and no training objective aimed
at it.

Whether suppression helps a classifier is a separate question. The suppressed images are computed from the radiograph and add no
information that a classifier could not extract itself, so a gain must come from the prior that the suppressor learned from CT or from
the way a network uses the extra images. A network that reads three images spends three times the computation; a gain can therefore be
credited to suppression only against a control of the same capacity. Seed variation and analytic flexibility can easily produce spurious
differences \cite{varoquaux2022failures,simmons2011falsepositive}, which is why every comparison here is paired, repeated over seeds and
fixed in advance.

CT can also teach a radiograph classifier directly. Knowledge distillation trains a student to match the softened outputs of a teacher
\cite{hinton2015distilling}, and cross-modal distillation transfers supervision between modalities through paired images
\cite{gupta2016crossmodal}. Each Qure.ai projection is built from one CT-RATE volume, so a CT classifier can teach an X-ray model on
images of the same patients, and one transformer family \cite{liu2022swinv2} can serve both modalities, which allows their heat maps to be
compared in the same patient. Whether such a transfer improves real-radiograph classification, aligns the explanations of the two models
or reduces hallucinated findings has not been tested.

We address these questions in three studies, each governed by an analysis plan that was pre-specified and hash-frozen before testing but
not deposited in a public registry. Study~I measures and reduces lesion erasure by CT-supervised bone suppressors, Study~II tests whether
classifiers gain from the suppressed images, and Study~III distils a Swin CT teacher into Swin radiograph classifiers. The main
contributions are:
\begin{itemize}
\item We define a counterfactual lesion audit that measures the fraction $A$ of an exactly known, additive lesion signal that a
 suppressor moves into its bone image, at sites shared by all models and reported beside suppression strength.
\item We introduce a twin lesion-invariance regulariser, trained on public CT-derived projections in a three-model ensemble with
 split-conformal intervals. On the locked VinDr-RibCXR set the ensemble absorbed 0.051 of an added lesion, against 0.149 for a released
 CT-supervised model and 0.178 for a matched control, at non-inferior bone-region PSNR.
\item We evaluate the suppressors on locked sets (paired dual-energy images, an annotated rib set, real nodules, downstream detection and
 classification) and externally on paediatric radiographs, with every hypothesis fixed in advance and every deviation disclosed.
\item We test bone-suppressed classifier inputs in five configurations, including a multi-stream network and a capacity-matched control
 (175 runs on three datasets). As separate streams the suppressed inputs did not raise macro-AUROC on any dataset (NIH $+0.0001$
 $[-0.0022, 0.0024]$); as extra channels they gave small gains (NIH up to $+0.006$).
\item We distil a 2.5D Swin CT teacher into Swin X-ray classifiers through paired projections and audit accuracy, CT--X-ray heat-map
 agreement, localisation and hallucination. Distillation did not improve NIH classification ($-0.0001$ $[-0.0016, 0.0014]$); Swin
 Grad-CAM++ located expert boxes less than half as often as DenseNet-121, yet noise flipped far fewer finding-free images for Swin
 (27--29\,\% against 78\,\%).
\end{itemize}

Fig.~\ref{fig:workflow} summarises the three studies. Section~\ref{sec:related} reviews related work and the gaps that motivate the
design. Section~\ref{sec:methods} formalises the audit and describes the data, the three studies and the statistical analysis.
Section~\ref{sec:results} reports the results of all pre-specified hypotheses, Section~\ref{sec:discussion} discusses them together with
the limitations, and Section~\ref{sec:conclusion} concludes.

\begin{figure*}[!t]
\centering
\safefig[0.9\textwidth]{figures/fig_workflow.pdf}
\caption{Workflow of the three studies, from data to models, evaluation and the main result. Study~I trains and audits CT-supervised
bone suppressors, Study~II tests the suppressed images as classifier inputs against a capacity-matched control, and Study~III distils a
2.5D Swin CT teacher into Swin X-ray classifiers. All hypotheses were fixed in analysis plans that were hash-frozen before testing;
training ran in cloud GPU kernels, and every reported number is computed by scripts from the saved outputs.}
\label{fig:workflow}
\end{figure*}
